\documentclass[10pt,a4paper]{amsart}
\usepackage[margin=19mm]{geometry}
\usepackage{amsmath,amssymb,mathtools}
\usepackage[scr=rsfs,cal=euler]{mathalfa}
\usepackage{xfrac}
\usepackage[unicode,hidelinks,bookmarks=false]{hyperref}
\newcommand{\ie}{i.e.\ }
\newcommand{\cf}{cf.\ }
\newcommand{\resp}{respectively\ }
\newcommand{\Subsection}{\subsection}
\providecommand{\nobreakdash}{\nobreak\hbox{-}}
\setlength{\emergencystretch}{2em}
\multlinegap0pt
\usepackage{amssymb}
\usepackage{bm}
\usepackage{mathtools}
\usepackage{xcolor}
\usepackage[shortlabels]{enumitem}
\newtheorem{proposition}{Proposition}
\newcommand{\E}{\mathbb{E}}
\title{Estimation of the sub-Gaussian parameter}
\author{Jason Liu, Min Xu, Jinchuan Xing}
\date{Source: arXiv:2606.06384v1 (CC BY 4.0)}
\begin{document}
\maketitle

\noindent\emph{Source and licence.} Estimation of the sub-Gaussian parameter. Version arXiv:2606.06384v1 (CC BY 4.0), distributed by arXiv under the Creative Commons Attribution 4.0 International licence, http://creativecommons.org/licenses/by/4.0/, as recorded in the arXiv licence metadata for this exact version and shown on the abstract page https://arxiv.org/abs/2606.06384v1. Full licence notice: \texttt{licenses/arxiv-2606.06384-CC-BY-4.0.txt}.\par\smallskip

\noindent\textit{Editorial scope.} Counted unit: the article's proposition prop:three-point-computation (source0.tex lines 1031--1040). The article's complete original proof of it is delivered with it (source0.tex lines 1041--1096, one \texttt{proof} environment of 56 lines). No mathematics is written, repaired or re-derived: the statement and the proof are the article's own lines, in the article's own notation, and no retained line is substituted or re-flowed.  Retained uncounted as the source-local context the proof needs: lines 1024--1027.  Nothing was omitted from any retained span, so the package carries no omission record.  No retained line contains a citation, so the wrapper carries no bibliography and no bibliography entry.  No retained line contains a figure environment, an image-inclusion command or drawing code, so no image is delivered and the package declares no asset dependency.  Editorial formatting only, in addition: an AMS \texttt{amsart} preamble that restates the article's own declarations and macros used by the retained lines, namely the environments \texttt{proposition} on separate counters, together with the standard packages those lines need.  The retained environments print in this document as Proposition 1 in order of appearance, whereas the article prints the same items with the numbering of its own full document.  The complete native source of the article as distributed in the arXiv e-print for this exact version is bundled unchanged as \texttt{sources/202184/source0.tex}.
\begin{proposition}
\label{prop:subgauss-scaling}
    Let $X \sim P$ be sub-Gaussian. If $aX \sim P_a$, then $\xi_*^2(P) = a^2\xi_*^2(P_a)$.
\end{proposition}

\begin{proposition}
    
\label{prop:three-point-computation}
    For a real sequence $(a_n)$, let $P := (1 - \frac1n)\delta_0 + \frac{1}{2n}\delta_{a_n} + \frac{1}{2n} \delta_{-a_n}$. The following hold:
    \begin{enumerate}
        \item[(a)] $\frac{a_n^2}{2\log 2n} \leq \xi_*^2(P) \leq \frac{a_n^2}{\log n}$. 
        %\item[(b)] Let $\lambda^*$ be a $[0, \infty)$-maximizer of $L(\lambda; P)$; we have $\lambda^* \leq 2a_n^{-1}\log 2n$. 
        \item[(b)] Put $a_n := (\log n)^\gamma$ for $\gamma \in [0, 1/2]$. Then $\delta_P(C) \leq  r_\gamma(C)$ for all $C > 0$.
    \end{enumerate}
\end{proposition}

\begin{proof}
    (a) Because we can rescale by $a_n$ and then apply Proposition~\ref{prop:subgauss-scaling}, it suffices to prove that if  $P = (1 - \frac1n)\delta_0 + \frac{1}{2n}\delta_{1} + \frac{1}{2n} \delta_{-1}$ then 
    \[
     \frac{1}{2 \log 2n}\leq \xi_*^2(P) \leq \frac{1}{\log n}.
    \]
    First we address the lower bound.  We have \begin{align*}
        M(2 \log 2n; P) &= \left(1 - \frac1n\right) + \frac{1}{2n}e^{2 \log 2n} + \frac{1}{2n}e^{-2 \log 2n} \\
        &= \left(1 - \frac1n\right) + \frac1n\left(\frac{(2n)^2 +(2n)^{-2}}{2}\right) \\
        &\geq \left(1 - \frac1n\right)  + 2n \geq 2n,
    \end{align*}
    so it follows that 
    \[
    \xi_*^2(P) \geq L(2 \log 2n; P) \geq \frac{2}{(2 \log 2n)^2}\log 2n \geq \frac{1}{2\log 2n}.
    \]
    For the upper bound, we will use the definition of $\xi_*^2$ and show that 
    \begin{equation}
        M(\lambda; P) \leq \exp\left(\frac{\lambda^2}{2\log n}\right).
    \label{eqn:mgf-upper-bound}
    \end{equation}
    There are two cases to handle. First, if $\lambda \leq \sqrt{2\log n}$, we have that 
    \[
    \log M(\lambda) \leq \frac{1}{n}\left(\frac{e^{\lambda} - e^{-\lambda}}{2} - 1\right) \leq\frac{1}{n}\left(e^{\lambda^2/2} - 1\right)
    \]
    by the inequalities $\log(1+x)\leq x$ and $\frac{e^x - e^{-x}}{2} \leq e^{x^2/2}$. Then since $s \mapsto (e^s - 1)/s$ is increasing and $\lambda^2/2 \leq \log n$ we have 
    \[
\frac{e^{\lambda^2/2}-1}{\lambda^2/2} \leq  \frac{e^{\log n} - 1}{\log n} \leq \frac{n}{\log n}
    \]
    so multiplying by $\lambda^2/2$ yields $e^{\lambda^2/2} - 1 \leq \frac{n\lambda^2}{2\log n}$. Therefore
    \[
    \log M(\lambda) \leq \frac{\lambda^2}{2 \log n},
    \]
    so $\lambda$ satisfies (\ref{eqn:mgf-upper-bound}) in this case. On the other hand, if $\lambda \geq \sqrt{2 \log n}$, then we have 
    \[
    \log M(\lambda) \leq \frac{1}{n}\left(\frac{e^{\lambda} - e^{-\lambda}}{2} - 1\right)  \leq \frac{e^{\lambda} - 1}{n}.
    \]
    By the AM-GM inequality, $\frac12(\log n + \frac{\lambda^2}{\log n})\geq \lambda$, so it follows that $e^{\lambda } \leq \exp(\frac12(\log n + \frac{\lambda^2}{\log n})) = \sqrt n \exp(\frac{\lambda^2}{2\log n})$. Combining this with the above display, it only remains to prove that $\sqrt n\exp(\frac{\lambda^2}{2\log n}) - 1  \leq n\exp(\frac{\lambda^2}{2\log n})$, but this holds trivially because $\frac{\lambda^2}{2\log n} > 1$. 
    \\
    %(b) Write $\lambda_n = 2a_n^{-1}\log 2n$. By a direct computation, we obtain that $L'(\lambda) > 0$ for all $\lambda \in (0, \lambda^*)$ and $L'(\lambda) < 0$ for all $\lambda \in (\lambda^*, \infty)$. It thus suffices to prove that $L'(\lambda_n) < 0$. Indeed, we have 
    %\[
    %L'(\lambda) = \frac{2}{\lambda^3}[\lambda\psi(\lambda) - 2 \psi(\lambda)],
    %\]
    %and by a direct computation, we have 
    %\[\frac{\lambda_nM'(\lambda_n)}{M(\lambda_n)} =  \lambda_na_n\frac{(2n - 1/8n^3)}{2n + 1 - 1/n + 1/8n^3} \leq \lambda_na_n =2\log 2n.
    %\]
    %Furthermore since $M(\lambda_n) \geq 2n$, $2\psi(\lambda_n) \geq 2\log 2n$. Putting these two inequalities together shows that $\lambda_n\psi(\lambda_n) < 2\psi(\lambda_n)$, and therefore $L'(\lambda_n) < 0$.
    %\\
    (b) Let $C > 0$ be arbitrary. If $2(\log n)^{1 - \gamma} \leq C$ then we have 
    \begin{align*}
        \delta_P(C) &\leq \sup_{|\lambda| \geq C} L(\lambda) \leq \sup_{|\lambda| \geq C} \frac{2}{\lambda^2} \log \E[e^{a_n\lambda}] \\ &\leq \sup_{|\lambda| \geq C}  \frac{2a_n}{\lambda} = \frac{2a_n}{C} \leq r_\gamma(C)
    \end{align*}
    where the last inequality follows by direct computation because $\gamma \geq 0$. 
    On the other hand, if $C \leq 2(\log n)^{1 - \gamma}$ then by (a), 
    \[
    \delta_P(C) \leq \xi_*^2(P) \leq (\log n)^{2\gamma - 1} = 2^{\frac{1 - 2\gamma}{1 - \gamma}}2^{\frac{2\gamma - 1}{1 - \gamma}} \frac{(\log n)^{\frac{2\gamma - 1}{1 - \gamma}}}{(\log n)^{\gamma\frac{2\gamma - 1}{1 - \gamma}}}=r_\gamma(2(\log n)^{1 - \gamma}) \leq r_\gamma(C).
    \]
\end{proof}

\end{document}
